\section{Semianelli termodinamici}
\label{sec:semirings}

Questa sezione richiama le definizioni di Marcolli e Thorngren (MT)~\citep{marcolli2014thermodynamic}.
Le Proposizioni~\ref{prop:shannon} e~\ref{prop:kl} sono note e le dimostriamo per completezza.
La Proposizione~\ref{prop:tsallis} è una derivazione nostra, che usa la distribuzione di Curado e Tsallis~\citep{curado1991generalized}.
Lavoriamo nelle coordinate min-plus, in cui la ``somma'' tropicale è il minimo e il ``prodotto'' è la somma ordinaria.
In queste coordinate un numero $x$ ha il ruolo di un'energia.

\subsection{Definizione}

Sia $\Delta_n=\{p\in[0,1]^n:\ \sum_ip_i=1\}$ il simplesso delle distribuzioni su $n$ elementi.
Sia $S$ una funzione di entropia su $\Delta_n$, e sia $T\ge0$ una temperatura.

\begin{definizione}[MT, Def.~4.1]
\label{def:thermo}
La somma termodinamica di $x_1,\dots,x_n\in\R\cup\{+\infty\}$ è
\begin{equation}
\bigoplus_{S,T}\,(x_1,\dots,x_n)=\min_{p\in\Delta_n}\Bigl[\sum_{i=1}^n p_ix_i-T\,S(p)\Bigr].
\label{eq:thermo-sum}
\end{equation}
Nel caso binario si scrive $x\oplus_{S,T}y=\min_{p\in[0,1]}[px+(1-p)y-TS(p)]$.
\end{definizione}

La struttura $(\R\cup\{+\infty\},\oplus_{S,T},+)$ è il semianello termodinamico.
Per $T=0$ la~\eqref{eq:thermo-sum} è il minimo, e si ottiene il semianello tropicale.
Il termine $\sum_ip_ix_i$ è un'energia media e il termine $-TS(p)$ è il costo entropico, quindi la~\eqref{eq:thermo-sum} è un'energia libera.

\begin{osservazione}[Proprietà algebriche, MT Teor.~3.1 e 4.2]
L'operazione $\oplus_{S,T}$ è commutativa se e solo se $S(p)=S(1-p)$.
L'elemento $+\infty$ è neutro se e solo se $S(0)=S(1)=0$.
L'operazione è associativa se e solo se $S$ soddisfa la relazione di ricorsività
\begin{equation}
S(p_1)+(1-p_1)\,S\!\Bigl(\frac{p_2}{1-p_1}\Bigr)=S(p_1+p_2)+(p_1+p_2)\,S\!\Bigl(\frac{p_1}{p_1+p_2}\Bigr).
\label{eq:assoc}
\end{equation}
A meno di una costante moltiplicativa, l'unica entropia che soddisfa le tre condizioni è quella di Shannon, $\Sh(p)=-p\log p-(1-p)\log(1-p)$.
\end{osservazione}

L'associatività conta per il message passing.
BP e SP aggregano molti messaggi in ingresso con la stessa operazione.
Se l'operazione non è associativa, il risultato dipende dall'ordine in cui si combinano i messaggi.
La legge distributiva generalizzata di Aji e McEliece~\citep{aji2000generalized} richiede un semianello commutativo, cioè una somma associativa e commutativa e un prodotto che distribuisce sulla somma.

\subsection{Shannon: la somma di Gibbs--Maslov}

\begin{proposizione}[MT, Prop.~4.3]
\label{prop:shannon}
Con l'entropia di Shannon $S=H(p)=-\sum_ip_i\log p_i$ e $T>0$ vale
\begin{equation}
\bigoplus_{\Sh,T}(x_1,\dots,x_n)=-T\log\sum_{i=1}^ne^{-x_i/T},
\label{eq:shannon-sum}
\end{equation}
e il minimo è raggiunto nella distribuzione di Gibbs $p_i^*=e^{-x_i/T}/\sum_je^{-x_j/T}$.
\end{proposizione}

\begin{proof}
Poniamo $Z=\sum_je^{-x_j/T}$ e $g_i=e^{-x_i/T}/Z$.
Allora $x_i=-T\log g_i-T\log Z$.
Sostituendo si ottiene, per ogni $p\in\Delta_n$,
\begin{equation*}
\sum_ip_ix_i-TH(p)=-T\log Z+T\sum_ip_i\log\frac{p_i}{g_i}=-T\log Z+T\,\KL(p\,\|\,g).
\end{equation*}
L'entropia relativa $\KL(p\|g)$ è non negativa e vale zero se e solo se $p=g$.
Quindi il minimo vale $-T\log Z$ ed è raggiunto in $p=g$.
\end{proof}

Questa identità è il principio variazionale di Gibbs.
La mappa $x\mapsto e^{-x/T}$ trasforma $\oplus_{\Sh,T}$ nella somma ordinaria e $+$ nel prodotto ordinario.
Quindi, a $T$ fisso, il semianello di Shannon è isomorfo al semianello $(\R_{\ge0},+,\times)$ di BP.
A $T$ fisso esso non contiene nuova algebra: $T$ è una scala.
L'informazione nuova sta nella famiglia in $T$ e nel limite $T\to0$, in cui la~\eqref{eq:shannon-sum} tende a $\min_ix_i$.
Il parametro $T$ coincide con il parametro $h$ della dequantizzazione di Maslov~\citep{litvinov2007maslov}.
Per questa ragione il semianello di Shannon si chiama anche semianello di Gibbs--Maslov.

\subsection{Entropia relativa}

\begin{proposizione}
\label{prop:kl}
Sia $\pi\in\Delta_n$ con $\pi_i>0$.
Con $S(p)=-\KL(p\|\pi)$ e $T>0$ vale
\begin{equation}
\bigoplus_{\KL(\cdot\|\pi),T}(x_1,\dots,x_n)=\min_{p\in\Delta_n}\Bigl[\sum_ip_ix_i+T\,\KL(p\,\|\,\pi)\Bigr]=-T\log\sum_{i=1}^n\pi_i\,e^{-x_i/T}.
\label{eq:kl-sum}
\end{equation}
\end{proposizione}

\begin{proof}
Si ripete la dimostrazione precedente con $g_i=\pi_ie^{-x_i/T}/Z_\pi$ e $Z_\pi=\sum_j\pi_je^{-x_j/T}$.
Si ottiene $\sum_ip_ix_i+T\KL(p\|\pi)=-T\log Z_\pi+T\KL(p\|g)$.
\end{proof}

MT studiano il caso binario nella Sezione~8.
Con $\pi$ uniforme si ha $\KL(p\|\pi)=\log n-H(p)$, e la~\eqref{eq:kl-sum} è la~\eqref{eq:shannon-sum} più la costante $T\log n$.
Con $\pi$ non uniforme l'operazione non è commutativa, perché il peso $\pi_i$ è legato alla posizione $i$.
La~\eqref{eq:kl-sum} è la forma che compare nelle equazioni di cavità: la distribuzione $\pi$ è la probabilità a priori delle configurazioni dei messaggi in ingresso, e $x_i$ è il costo in energia della configurazione $i$ (Sezione~\ref{sec:spy}).

\subsection{Rényi}

L'entropia di Rényi di ordine $\alpha>0$, $\alpha\ne1$, è
\begin{equation}
\Ry_\alpha(p)=\frac{1}{1-\alpha}\log\sum_ip_i^{\alpha},
\label{eq:renyi}
\end{equation}
e tende a $H(p)$ per $\alpha\to1$~\citep{renyi1961measures}.
È additiva su sistemi indipendenti ma non soddisfa la~\eqref{eq:assoc}.
Quindi $\oplus_{\Ry_\alpha,T}$ è commutativa ma non associativa.
Il Lemma~6.1 di MT mostra che il difetto di associatività si corregge con la permutazione $(p_1,p_2,p_3)\mapsto(p_3,p_2,p_1)$.
Per $T\to0$ anche $\oplus_{\Ry_\alpha,T}$ tende al minimo.
Nella Sezione~\ref{sec:qaxis} l'entropia di Rényi non entra come operazione di messaggio ma come misura dei pesi dei cluster.

\subsection{Tsallis e la somma deformata}

L'entropia di Tsallis di indice $q>0$, $q\ne1$, è
\begin{equation}
S_q(p)=\frac{1-\sum_ip_i^{\,q}}{q-1},
\label{eq:tsallis}
\end{equation}
e tende a $H(p)$ per $q\to1$~\citep{tsallis1988possible}.
Non è additiva su sistemi indipendenti: $S_q(A\star B)=S_q(A)+S_q(B)+(1-q)S_q(A)S_q(B)$.
Anche $\oplus_{S_q,T}$ non è associativa.
MT (Sezione~7.1, eq.~(7.4)) introducono un'operazione deformata.
Nelle coordinate min-plus essa si scrive
\begin{equation}
x\oplus_{S,T,q}y=\min_{p\in[0,1]}\bigl[p^{q}x+(1-p)^{q}y-T\,S(p)\bigr].
\label{eq:tsallis-op}
\end{equation}
Il Teorema~7.2 di MT afferma che l'operazione~\eqref{eq:tsallis-op} è associativa se e solo se $S$ soddisfa la versione deformata della~\eqref{eq:assoc}, in cui i fattori $(1-p_1)$ e $(p_1+p_2)$ sono elevati alla $q$.
Il Teorema~7.1 afferma che l'entropia di Tsallis è l'unica entropia commutativa, con la proprietà di identità, che soddisfa questa condizione.
La deformazione sposta la potenza $q$ dal termine di entropia al termine di energia: l'energia media diventa $U_q=\sum_ip_i^{q}x_i$.
Questa è la ``seconda scelta'' di Curado e Tsallis~\citep{curado1991generalized}.

\begin{proposizione}[Forma chiusa della somma di Tsallis]
\label{prop:tsallis}
Siano $x_i\ge0$ e $T>0$.
Poniamo $\lnq u=(u^{1-q}-1)/(1-q)$ ed $\eq(u)=[1+(1-q)u]_+^{1/(1-q)}$, dove $[z]_+=\max(z,0)$.
Allora
\begin{equation}
F_q(x;T)\equiv\min_{p\in\Delta_n}\Bigl[\sum_ip_i^{\,q}x_i-T\,S_q(p)\Bigr]=-T\,\lnq\sum_{i=1}^n\eq(-x_i/T),
\label{eq:tsallis-closed}
\end{equation}
e il minimo è raggiunto in $p_i^*\propto\eq(-x_i/T)$.
Per $q\to1$ la~\eqref{eq:tsallis-closed} tende alla~\eqref{eq:shannon-sum}.
\end{proposizione}

\begin{proof}
Riscriviamo la funzione obiettivo come
\begin{equation*}
\mathcal F(p)=\sum_ip_i^{\,q}x_i+\frac{T}{q-1}\Bigl(\sum_ip_i^{\,q}-1\Bigr)=\frac{T}{q-1}\Bigl[\sum_ip_i^{\,q}\,u_i-1\Bigr],\qquad u_i=1+\frac{(q-1)x_i}{T}.
\end{equation*}
Caso $q>1$.
Si ha $u_i>0$, e $\mathcal F$ è strettamente convessa, perché $p\mapsto p^q$ è strettamente convessa e il coefficiente $u_iT/(q-1)$ è positivo.
Il punto stazionario del lagrangiano $\mathcal F(p)+\lambda(\sum_ip_i-1)$ è quindi il minimo globale.
La condizione di stazionarietà è $q\,p_i^{\,q-1}u_iT/(q-1)=-\lambda$, quindi $p_i\propto u_i^{-1/(q-1)}=\eq(-x_i/T)$.
Poniamo $Z_q=\sum_ju_j^{-1/(q-1)}$, così che $p_i=u_i^{-1/(q-1)}/Z_q$.
Allora $p_i^{\,q}u_i=u_i^{-1/(q-1)}Z_q^{-q}$, e $\sum_ip_i^{\,q}u_i=Z_q^{1-q}$.
Il valore minimo è $T(Z_q^{1-q}-1)/(q-1)=-T\lnq Z_q$.

Caso $0<q<1$.
Si ha $\mathcal F(p)=\frac{T}{1-q}\bigl[1-\sum_ip_i^{\,q}u_i\bigr]$.
Se $u_i\le0$, il termine $-p_i^{\,q}u_i$ è non negativo e il minimo pone $p_i=0$.
Questo è il taglio (cutoff) della distribuzione di Tsallis per $q<1$.
Sugli indici con $u_i>0$ la funzione $-\sum_ip_i^qu_i$ è strettamente convessa, e il calcolo del caso precedente dà lo stesso risultato con $Z_q=\sum_{u_j>0}u_j^{1/(1-q)}=\sum_j\eq(-x_j/T)$.

Per $q\to1$ si ha $\eq(u)\to e^{u}$ e $\lnq\to\log$.
\end{proof}

Abbiamo verificato numericamente la~\eqref{eq:tsallis-closed} con una minimizzazione diretta sul simplesso per alcuni valori di $(q,T)$.

\begin{corollario}[Limite $T\to0$ della somma di Tsallis]
\label{cor:tsallis-T0}
Siano $x_i>0$.
Per $0<q\le1$ si ha $\lim_{T\to0}F_q(x;T)=\min_ix_i$.
Per $q>1$ si ha
\begin{equation}
\lim_{T\to0}F_q(x;T)=\min_{p\in\Delta_n}\sum_ip_i^{\,q}x_i=\Bigl(\sum_ix_i^{-1/(q-1)}\Bigr)^{-(q-1)}.
\label{eq:power-mean}
\end{equation}
Questo valore è strettamente minore di $\min_ix_i$ se almeno due $x_i$ sono finiti, e tende a $\min_ix_i$ per $q\to1^+$.
\end{corollario}

\begin{proof}
Per $T\to0$ il termine $-TS_q$ scompare, perché $S_q$ è limitata sul simplesso.
Resta il minimo di $U_q(p)=\sum_ip_i^qx_i$.
Per $q<1$ la funzione $U_q$ è concava, quindi il minimo è in un vertice del simplesso e vale $\min_ix_i$.
Per $q=1$ la funzione è lineare e il risultato è lo stesso.
Per $q>1$ la funzione è strettamente convessa.
La stazionarietà dà $p_i\propto x_i^{-1/(q-1)}$, e sostituendo si ottiene la~\eqref{eq:power-mean}.
Ogni termine della somma in~\eqref{eq:power-mean} è positivo, quindi la somma supera $(\min_ix_i)^{-1/(q-1)}$ e il valore è minore di $\min_ix_i$.
Per $q\to1^+$ l'esponente $1/(q-1)$ tende a infinito, e la media di potenze tende al minimo.
\end{proof}

Il corollario dice che la somma di Tsallis deformata con $q>1$ non si dequantizza al semianello tropicale.
Il suo limite $T\to0$ è una media di potenze, non il minimo.
I due limiti iterati danno entrambi il minimo: $\lim_{T\to0}\lim_{q\to1}F_q=\min_ix_i$ e $\lim_{q\to1^+}\lim_{T\to0}F_q=\min_ix_i$.
Il quadrato dei limiti quindi commuta, ma a $q>1$ fisso il limite $T\to0$ ha un vertice diverso dal minimo.
Questo fatto compare nel diagramma della Sezione~\ref{sec:diagrams}.

\subsection{Funzione successore e cumulanti}

MT definiscono la funzione successore $\lambda(x,T)=x\oplus_{S,T}0$ (Sezione~9).
Essa è la trasformata di Legendre di $TS$.
La Proposizione~9.4 di MT afferma che $-\lambda/T$ è la funzione generatrice dei cumulanti dell'energia.
Nel caso di Shannon questa è la relazione usuale fra energia libera e cumulanti.
Nella Sezione~\ref{sec:spy} la useremo nella forma seguente: $G(y)=-y\Phi(y)$ è la funzione generatrice dei cumulanti dell'energia degli stati, per variabile, e la trasformata~\eqref{eq:Phi-y} è la relazione di Gärtner--Ellis fra questa funzione e la complessità~\citep{touchette2009large}.
